% part: intuitionistic-logic
% chapter: propositions-as-types
% section: types

\documentclass[../../../include/open-logic-section]{subfiles}

\begin{document}

\olfileid{int}{pty}{typ}

\olsection{类型}

像 $(MN)$ 这样的证明项，其本身并不能唯一确定它所对应的证明。这是因为自由变元（即公理的证明项）并未指明对应公理中的公式。然而，如果我们在一个语境~$\Gamma$ 中指明这些公式，那么这个证明\emph{就}能被唯一确定。

\begin{defn}
证明项~$N$ 的一个\emph{语境}是一个集合 $\Gamma$，它满足：对于~$N$ 的每个自由变元，$\Gamma$ 都恰好包含一个对 $x : !A$。
\end{defn}

注意，该定义并不要求 $\Gamma$ 包含的 $x:!A$ \emph{仅}针对 $N$ 的自由变元。例如，$\Gamma =
\{x:!A, y:!B, z:!C\}$ 就是 $\pair{x}{y}$ 的一个语境。

\begin{defn}\ollabel{defn:type}
  设 $\Gamma$ 是~$N$ 的一个语境。$N$ 的\emph{类型}（相对于~$\Gamma$）归纳定义如下：
  \begin{enumerate}
  \item $x$ 的类型是 $!A$，当且仅当 $x:!A \in \Gamma$。
  \item 若 $N$ 相对于 $x:!A, \Gamma$ 有类型~$!B$，则 $\lambd[\typeof{x}{!A}][N]$ 的类型为 $!A \lif !B$。
  \item 若 $N$ 有类型 $!A \lif !B$ 且 $M$ 有类型~$!A$，则 $(NM)$ 有类型~$!B$。
  \item 若 $N$ 有类型 $!A$ 且 $M$ 有类型~$!B$，则 $\pair{N}{M}$ 有类型 $!A \land !B$。
  \item 若 $N$ 有类型~$!A_1 \land !A_2$，则 $\proj{i}{N}$ 有类型~$!A_i$。
  \item 若 $N$ 有类型 $!A$，则
    $\inj[!B]{i}{!M}$ 在 $i=1$ 时有类型 $!A \lor !B$，在 $i=2$ 时有类型 $!B \lor !A$。
  \item 若 $M$ 有类型 $!A \lor !B$，$N_1$ 相对于 $x_1:!A, \Gamma$ 有类型 $!C$，且 $N_2$ 相对于 $x_2:!B, \Gamma$ 有类型~$!C$，则 $\dcase{M}{x_1}{N_1}{x_2}{N_2}$ 有类型~$!C$。
  \item 若 $N$ 有类型~$\lfalse$，则 $\abort{!A}{N}$ 有类型~$!A$。
  \end{enumerate}
\end{defn}

\begin{prop}
若一个证明项~$N$ 有类型（即它是一个正确的证明项），则它相对于它的一个语境~$\Gamma$ 有且仅有一个类型。
\end{prop}

\begin{proof}
  对~$N$ 归纳证明。
\end{proof}

当 $N$ 相对于 $\Gamma$ 的类型是~$!A$ 时，我们写作 $\Gamma \Sequent N : !A$。\olref{defn:type}中的分款看起来应该很熟悉：它们正是 \Log{tN2i} 的规则，只有一点细微差别：语境~$\Gamma$ 在整个过程中保持不变。我们可以将这个定义编码为一个演算系统 \Log{tN3i}，其规则如 \olref{tab:tN2ip} 所示。

\subfile{rules-tN3}

证明项是某个版本的$\lambd$-演算中的项，即\emph{带积、和与空类型的带类型$\lambd$-演算}。经典的$\lambd$-演算只包含由变元、\emph{$\lambd$-抽象}~$\lambd[x][N]$ 和作用 $(MN)$ 构建的项，并且缺少带类型的$\lambd$-抽象 $\lambd[\typeof{x}{!A}][N]$ 中对~$!A$ 的标注。$\lambd$-演算是一种图灵完备的计算模型（即每个部分可计算函数都可以在其中用一个项来表示）。由证明项和类型规则给出的这个$\lambd$-演算版本是一个受限的计算模型，但基本思想相同。它的限制在于类型指派：只有\emph{良类型的}项，即相对于某个语境~$\Gamma$ 有类型的证明项，才是合法的项。例如，$\lambd[x][(xx)]$ 是经典无类型$\lambd$-演算中的一个合法项，但对于任何 $!A$，$\lambd[\typeof{x}{!A}][(xx)]$ 都不是良类型的。这是因为类型规则要求，为了使 $(xx)$ 有类型，第一个~$x$ 必须具有类型 $!A \lif !B$ 而第二个 $x$ 具有类型~$!A$，而这是不可能的，因为 $!A$ 是 $!A \lif !B$ 的一个真子公式。

我们可以直观地将类型理解为：具有该类型的项所命名的对象的类别。因此，一个类型为 $!A \land !B$ 的项挑选的是一对元素，其第一个分量的类型是~$!A$，第二个分量的类型是~$!B$，即积~$!A \times !B$ 中的一个元素。一个类型为 $!A \lif !B$ 的项是一个从类型 $!A$ 的对象到类型~$!B$ 的对象的函数。一个类型为 $!A \lor !B$ 的项要么是类型 $!A$ 的对象，要么是类型~$!B$ 的对象，并且还附带说明它是哪一个的信息。这类似于两个集合 $!A$ 和 $!B$ 的无交并，定义为 $\{\tuple{i,x} : i = 1 \text{ 或 } 2, x \in !A \text{ 若 } i = 1, x \in !B \text{ 若 } i=2\}$。类型 $\lfalse$ 是空集，也就是说，如果一个项具有该类型，它就不能命名任何东西。由于自然演绎（包括 \Log{tN3i}）是一致的，除非存在未解除的假设，否则不能推导出~$\lfalse$，因此不存在类型为~$\lfalse$ 的闭证明项（不含自由变元的项）。

带类型的 $\lambd$-演算的算子会产生一些项，这些项直观地命名了正确类别的对象，前提是它所作用的那些项命名了某些对象。例如，“若 $N_1 : !A$ 且 $N_2 : !B$，则 $\pair{N_1}{N_2} : !A \land !B$” 的意思是：如果 $N_1$ 命名了一个类型为~$!A$ 的对象，并且 $N_2$ 命名了一个类型为~$!B$ 的对象，那么 $\pair{N_1}{N_2}$ 就命名了一个类型为 $!A \land !B$ 的对象。

\end{document}